\section{Iteration chunking via \code{iteration\_precision}}
\label{sec:chunking}

A few of the base kernels described in \cref{sec:mandel-base-kernels}, which exclude linear approximation or
perturbation, are templated on an integer \code{iteration\_precision}
(\(1,2,4,8,16\)) and unroll multiple Mandelbrot steps inside the loop:

\begin{itemize}
  \item Each loop iteration performs \code{iteration\_precision} updates.
  \item The counter \code{iter} increases by that amount.
  \item The maximum iteration \code{n\_iterations} is adjusted so \code{iter}
        does not exceed the requested limit.
\end{itemize}

This approach reduces loop overhead. The trade-off is that the escape predicate
is typically checked only once per chunk; therefore the reported escape
iteration can be larger than the true first-escape iteration by up to
\code{iteration\_precision - 1}. For strict escape iteration counts (e.g., for
continuous/smooth coloring based on the first bailout), use a chunk size of 1 or
insert bailout checks within the unrolled body.

\paragraph{Worked chunk example.}
For $c=1$, the exact orbit begins $0,1,2,5,26$. The first value strictly outside
the radius-two circle is $z_3=5$. A loop that checks only after four updates first
observes escape at $z_4=26$, one update late. The example counts recurrence updates;
an individual kernel's stored-counter convention must also be considered when
comparing output counts. More generally, a chunk of $k$ updates can postpone the
observation by at most $k-1$ updates, provided arithmetic remains usable until the
check. Large escaped values may overflow before the end of a chunk. The iteration
budget must also be handled explicitly when it is not a multiple of $k$.

This optimization was mainly explored for educational purposes; in practice, the
benefit is small compared to other optimizations such as perturbation (\cref{sec:perturbation-concept}) and linear
approximation (\cref{sec:approx-accel}).

The kernels described above operate on standard IEEE~754 types.  The next
chapter introduces the extended-precision and high-dynamic-range numeric types
that \FractalShark{} employs when native precision is insufficient.
